% ======================================================================
% Chapter 8: Invitation: Ultraviolet Cutoffs and Critical Fluctuations
% (book pages 265-274; PDF pages 288-297)
% ======================================================================
\chapter{Invitation: Ultraviolet Cutoffs and Critical Fluctuations}

The main purpose of Part~II of this book is to develop a general theory of
renormalization.  This theory will explain the origin of ultraviolet
divergences in field theory and will indicate when these divergences can be
removed systematically.  It will also give a way to convert the divergences of
Feynman diagrams from a problem into a tool.  We will apply this tool to
study the asymptotic large- or small-momentum behavior of field theory
amplitudes.

When we first encountered an ultraviolet divergence in the calculation of the
one-loop vertex correction in Section~6.3, it seemed an aberration that ought
to disappear before it caused us too much discomfort.  In Chapter~7 we saw
further examples of ultraviolet-divergent diagrams, enough to convince us that
such divergences occur ubiquitously in Feynman diagram computations.  Thus it
is necessary for anyone studying field theory to develop a point of view
toward these divergences.  Most people begin with the belief that any theory
that contains divergences must be nonsense.  But this viewpoint is overly
restrictive, since it excludes not only quantum field theory but even the
classical electrodynamics of point particles.

With some experience, one might adopt a more permissive attitude of peaceful
coexistence with the divergences: One can accept a theory with divergences, as
long as they do not appear in physical predictions.  In Chapter~7 we saw that
all of the divergences that appear in the one-loop radiative corrections to
electron scattering from a heavy target can be eliminated by consistently
eliminating the bare values of the mass and charge of the electron in favor of
their measured physical values.  In Chapter~10, we will argue that all of the
ultraviolet divergences of QED, in all orders of perturbation theory, can be
eliminated in this way.  Thus, as long as one is willing to consider the mass
and charge of the electron as measured parameters, the predictions of QED
perturbation theory will always be free of divergences.  We will also show in
Chapter~10 that QED belongs to a well-defined class of field theories in which
all ultraviolet divergences are removed after a fixed small number of physical
parameters are taken from experiment.  These theories, called
\emph{renormalizable} quantum field theories, are the only ones in which
perturbation theory gives well-defined predictions.

Ideally, though, one should take the further step of trying to understand
physically why the divergences appear and why their effects are more severe in
some theories than in others.  This direct approach to the divergence problem
was pioneered in the 1960s by Kenneth Wilson.  The crucial insights needed to
solve this problem emerged from a correspondence, discovered by Wilson and
others, between quantum field theory and the statistical physics of magnets
and fluids.  Wilson's approach to renormalization is the subject of
Chapter~12.  The present chapter gives a brief introduction to the issues in
condensed matter physics that have provided insight into the problem of
ultraviolet divergences.

\subsection*{Formal and Physical Cutoffs}

Ultraviolet divergences signal that quantities calculated in a quantum field
theory depend on some very large momentum scale, the ultraviolet cutoff.
Equivalently, in position space, divergent quantities depend on some very
small distance scale.

The idea of a small-distance cutoff in the continuum description of a system
occurs in classical field theories as well.  Typically the cutoff is at the
scale of atomic distances, where the continuum description no longer applies.
However, the size of the cutoff manifests itself in certain parameters of the
continuum theory.  In fluid dynamics, for instance, parameters such as the
viscosity and the speed of sound are of just the size one would expect by
combining typical atomic radii and velocities.  Similarly, in a magnet, the
magnetic susceptibility can be estimated by assuming that the energy cost of
flipping an electron spin is on the order of a tenth of an eV, as we would
expect from atomic physics.  Each of these systems possesses a natural
ultraviolet cutoff at the scale of an atom; by understanding the physics at
the atomic scale, we can compute the parameters that determine the physics on
larger scales.

In quantum field theory, however, we have no precise knowledge of the
fundamental physics at very short distance scales.  Thus, we can only measure
parameters such as the physical charge and mass of the electron, not compute
them from first principles.  The presence of ultraviolet divergences in the
relations between these physical parameters and their bare values is a sign
that these parameters are controlled by the unknown short-distance physics.

Whether we know the fundamental physics at small distance scales or not, we
need two kinds of information in order to write an effective theory for
large-distance phenomena.  First, we must know how many parameters from the
small distance scale are relevant to large-distance physics.  Second, and more
importantly, we must know what degrees of freedom from the underlying theory
appear at large distances.

In fluid mechanics, it is something of a miracle, from the atomic point of
view, that any large-distance degrees of freedom even exist.  Nevertheless, the
equations that express the transport of energy and mass over large distances
do have smooth, coherent solutions.  The large-distance degrees of freedom are
the flows that transport these conserved quantities, and sound waves of long
wavelength.

In quantum field theory, the large-distance physics involves only those
particles that have masses that are very small compared to the fundamental
cutoff scale.  These particles and their dynamics are the quantum analogues of
the large-scale flows in fluid mechanics.  The simplest way to naturally
arrange for such particles to appear is to make use of particles that
naturally have zero mass.  So far in this book, we have encountered two types
of particles whose mass is precisely zero, the photon and the chiral fermion.
(In Chapter~11 we will meet one further naturally massless particle, the
Goldstone boson.)  We might argue that QED exists as a theory on scales much
larger than its cutoff because the photon is naturally massless and because
the left- and right-handed electrons are very close to being chiral fermions.

There is another way that particles of zero or almost zero mass can arise in
quantum field theory: We can simply tune the parameters of a scalar field
theory so that the scalar particles have masses small compared to the cutoff.
This method of introducing particles with small mass seems arbitrary and
unnatural.  Nevertheless, it has an analogue in statistical mechanics that is
genuinely interesting in that discipline and can teach us some important
lessons.

Normally, in a condensed matter system, the thermal fluctuations are
correlated only over atomic distances.  Under special circumstances, however,
they can have much longer range.  The clearest example of this phenomenon
occurs in a ferromagnet.  At high temperature, the electron spins in a magnet
are disorganized and fluctuating; but at low temperature, these spins align to
a fixed direction.%
\footnote{In a real ferromagnet, the long-range magnetic dipole-dipole
interaction causes the state of uniform magnetization to break up into an
array of magnetic domains.  In this book, we will ignore this interaction and
think of a magnetic spin as a pure orientation.  It is this idealized system
that is directly analogous to a quantum field theory.}  Let us think about how
this alignment builds up as the temperature of the magnet is lowered.  As the
magnet cools from high temperature, clusters of correlated spins become larger
and larger.  At a certain point---the temperature of magnetization---the entire
sample becomes a single large cluster with a well-defined macroscopic
orientation.  Just above this temperature, the magnet contains large clusters
of spins with a common orientation, which in turn belong to still larger
clusters, such that the orientations on the very largest scale are still
randomized through the sample.  This situation is illustrated in
Fig.~\ref{fig:8.1}.  Similar behavior occurs in the vicinity of any other
second-order phase transition, for example, the order-disorder transition in
binary alloys, the critical point in fluids, or the superfluid transition in
Helium-4.

\begin{figure}[t]
  \centering
  \includegraphics[width=0.55\textwidth]{fig8.1.png}
  \caption{Clusters of oriented spins near the critical point of a
    ferromagnet.}
  \label{fig:8.1}
\end{figure}

The natural description of these very long wavelength fluctuations is in terms
of a fluctuating continuum field.  At the lowest intuitive level, we might
substitute quantum for statistical fluctuations and try to describe this
system as a quantum field theory.  In Section~9.3 we will derive a somewhat
more subtle relation that makes a precise connection between the statistical
and the quantum systems.  Through this connection, the behavior of any
statistical system near a second-order phase transition can be translated into
the behavior of a particular quantum field theory.  This quantum field theory
has a field with a mass that is very small compared to the basic atomic scale
and that goes to zero precisely at the phase transition.

But this connection seems to compound the problem of ultraviolet divergences
in quantum field theory: If the wealth of phase transitions observed in Nature
generates a similar wealth of quantum field theories, how can we possibly
define a quantum field theory without detailed reference to its origins in
physics at the scale of its ultraviolet cutoff?  Saying that a quantum field
theory makes predictions independent of the cutoff would be equivalent to
saying that the statistical fluctuations in the neighborhood of a critical
point are independent of whether the system is a magnet, a fluid, or an alloy.
But is this statement so obviously incorrect?  By reversing the logic, we
would find that quantum field theory makes a remarkably powerful prediction
for condensed matter systems, a prediction of \emph{universality} for the
statistical fluctuations near a critical point.  In fact, this prediction is
verified experimentally.

A major theme of Part~II of this book will be that these two ideas---cutoff
independence in quantum field theory and universality in the theory of
critical phenomena---are naturally the same idea, and that understanding
either of these ideas gives insight into the other.

\subsection*{Landau Theory of Phase Transitions}

To obtain a first notion of what could be universal in the phenomena of phase
transitions, let us examine the simplest continuum theory of second-order
phase transitions, due to Landau.

First we should review a little thermodynamics and clarify our nomenclature.
In thermodynamics, a \emph{first-order phase transition} is a point across
which some thermodynamic variable (the density of a fluid, or the
magnetization of a ferromagnet) changes discontinuously.  At a phase
transition point, two quite distinct thermodynamic states (liquid and gas, or
magnetization parallel and antiparallel to a given axis) are in equilibrium.
The thermodynamic quantity that changes discontinuously across the
transition, and that characterizes the difference of the two competing phases,
is called the \emph{order parameter}.  In most circumstances, it is possible
to change a second thermodynamic parameter in such a way that the two
competing states move closer together in the thermodynamic space, so that at
some value of this parameter, these two states become identical and the
discontinuity in the order parameter disappears.  This endpoint of the line of
first-order transitions is called a \emph{second-order phase transition}, or,
more properly, a \emph{critical point}.  Viewed from the other direction, a
critical point is a point at which a single thermodynamic state bifurcates
into two macroscopically distinct states.  It is this bifurcation that leads
to the long-ranged thermal fluctuations discussed in the previous section.

A concrete example of this behavior is exhibited by a ferromagnet.  Let us
assume for simplicity that the material we are discussing has a preferred axis
of magnetization, so that at low temperature, the system will have its spins
ordered either parallel or antiparallel to this axis.  The total magnetization
along this axis, $M$, is the order parameter.  At low temperature, application
of an external magnetic field $H$ will favor one or the other of the two
possible states.  At $H = 0$, the two states will be in equilibrium; if $H$ is
changed from a small negative to a small positive value, the thermodynamic
state and the value of $M$ will change discontinuously.  Thus, for any fixed
(low) temperature, there is a first-order transition at $H = 0$.  Now consider
the effect of raising the temperature: The fluctuation of the spins increases
and the value of $|M|$ decreases.  At some temperature $T_c$ the system ceases
to be magnetized at $H = 0$.  At this point, the first-order phase transition
disappears and the two competing thermodynamic states coalesce.  The system
thus has a critical point at $T = T_c$.  The location of these various
transitions in the $H$--$T$ plane is shown in Fig.~\ref{fig:8.2}.

\begin{figure}[t]
  \centering
  \includegraphics[width=0.5\textwidth]{fig8.2.png}
  \caption{Phase diagram in the $H$--$T$ plane for a uniaxial ferromagnet.}
  \label{fig:8.2}
\end{figure}

Landau described this behavior by the use of the Gibbs free energy $G$; this
is the thermodynamic potential that depends on $M$ and $T$, such that
\begin{equation}
  H = \frac{\partial G}{\partial M}\Big|_{T} .
  \label{eq:8.1}
\end{equation}
He suggested that we concentrate our attention on the region of the critical
point: $T\approx T_c$, $M\approx 0$.  Then it is reasonable to expand $G(M)$
as a Taylor series in $M$.  For $H = 0$, we can write
\begin{equation}
  G(M) = A(T) + B(T)M^2 + C(T)M^4 + \cdots .
  \label{eq:8.2}
\end{equation}
Because the system has a symmetry under $M\to -M$, $G(M)$ can contain only
even powers of $M$.  Since $M$ is small, we will ignore the higher terms in
the expansion.  Given Eq.~\eqref{eq:8.2}, we can find the possible values of
$M$ at $H = 0$ by solving
\begin{equation}
  0 = \frac{\partial G}{\partial M} = 2B(T)M + 4C(T)M^3 .
  \label{eq:8.3}
\end{equation}
If $B$ and $C$ are positive, the only solution is $M = 0$.  However, if
$C > 0$ but $B$ is negative below some temperature $T_c$, we have a
nontrivial solution for $T < T_c$, as shown in Fig.~\ref{fig:8.3}.  More
concretely, approximate for $T\approx T_c$:
\begin{equation}
  B(T) = b(T-T_c),\qquad C(T) = c .
  \label{eq:8.4}
\end{equation}
Then the solution to Eq.~\eqref{eq:8.3} is
\begin{equation}
  M \sim \begin{cases}
    0 & \text{for } T > T_c;\\
    \pm \bigl\{ (b/2c)(T_c - T) \bigr\}^{1/2} & \text{for } T < T_c .
  \end{cases}
  \label{eq:8.5}
\end{equation}
This is just the qualitative behavior that we expect at a critical point.

\begin{figure}[t]
  \centering
  \includegraphics[width=0.5\textwidth]{fig8.3.png}
  \caption{Behavior of the Gibbs free energy $G(M)$ in Landau theory, at
    temperatures above and below the critical temperature.}
  \label{fig:8.3}
\end{figure}

To find the value of $M$ at nonzero external field, we could solve
Eq.~\eqref{eq:8.1} with the left-hand side given by \eqref{eq:8.2}.  An
equivalent procedure is to minimize a new function, related to
\eqref{eq:8.2}.  Define
\begin{equation}
  G(M, H) = A(T) + B(T)M^2 + C(T)M^4 - HM .
  \label{eq:8.6}
\end{equation}
Then the minimum of $G(M,H)$ with respect to $M$ at fixed $H$ gives the value
of $M$ that satisfies Eq.~\eqref{eq:8.1}.  The minimum is unique except when
$H = 0$ and $T < T_c$, where we find the double minimum in the second line of
\eqref{eq:8.5}.  This is consistent with the phase diagram shown in
Fig.~\ref{fig:8.2}.

To study correlations in the vicinity of the phase transition, Landau
generalized this description further by considering the magnetization $M$ to
be the integral of a local spin density:
\begin{equation}
  M = \int d^3x\ s(\mathbf{x}) .
  \label{eq:8.7}
\end{equation}
Then the Gibbs free energy \eqref{eq:8.6} becomes the integral of a local
function of $s(\mathbf{x})$,
\begin{equation}
  G = \int d^3x\ \Bigl[ \tfrac{1}{2}(\nabla s)^2 + b(T-T_c)s^2 + cs^4 - Hs
    \Bigr],
  \label{eq:8.8}
\end{equation}
which must be minimized with respect to the field configuration
$s(\mathbf{x})$.  The first term is the simplest possible way to introduce the
tendency of nearby spins to align with one another.  We have rescaled
$s(\mathbf{x})$ so that the coefficient of this term is set to $1/2$.  In
writing this free energy integral, we could even consider $H$ to vary as a
function of position.  In fact, it is useful to do that; we can turn on
$H(\mathbf{x})$ near $\mathbf{x} = 0$ and see what response we find at another
point.

The minimum of the free energy expression \eqref{eq:8.8} with respect to
$s(\mathbf{x})$ is given by the solution to the variational equation
\begin{equation}
  0 = \frac{\delta G[s]}{\delta s(\mathbf{x})} = -\nabla^2 s + 2b(T-T_c)s
    + 4cs^3 - H(\mathbf{x}) .
  \label{eq:8.9}
\end{equation}
For $T > T_c$, where the macroscopic magnetization vanishes and so
$s(\mathbf{x})$ should be small, we can find the qualitative behavior by
ignoring the $s^3$ term.  Then $s(\mathbf{x})$ obeys a linear equation,
\begin{equation}
  \bigl( -\nabla^2 + 2b(T-T_c) \bigr)\,s(\mathbf{x}) = H(\mathbf{x}) .
  \label{eq:8.10}
\end{equation}
To study correlations of spins, we will set
\begin{equation}
  H(\mathbf{x}) = H_0\,\delta^{(3)}(\mathbf{x}) .
  \label{eq:8.11}
\end{equation}
The resulting configuration $s(\mathbf{x})$ is then the Green's function of
the differential operator in Eq.~\eqref{eq:8.10}, so we call it $D(\mathbf{x})$:
\begin{equation}
  \bigl( -\nabla^2 + 2b(T-T_c) \bigr)\,D(\mathbf{x}) = H_0\,\delta^{(3)}(\mathbf{x}).
  \label{eq:8.12}
\end{equation}
This Green's function tells us the response at $\mathbf{x}$ when the spin at
$\mathbf{x} = 0$ is forced into alignment with $H$.  In Sections~9.2 and~9.3
we will see that $D(\mathbf{x})$ is also proportional to the zero-field
spin-spin correlation function in the thermal ensemble,
\begin{equation}
  D(\mathbf{x}) \propto \langle s(\mathbf{x})s(0)\rangle
  = \sum_{\text{all } s(\mathbf{x})} s(\mathbf{x})\,s(0)\,
    e^{-\mathcal{H}/k_BT} ,
  \label{eq:8.13}
\end{equation}
where $\mathcal{H}$ is the Hamiltonian of the magnetic system.

The solution to Eq.~\eqref{eq:8.12} can be found by Fourier transformation:
\begin{equation}
  D(\mathbf{x}) = \int \frac{d^3k}{(2\pi)^3}\
    \frac{H_0\,e^{i\mathbf{k}\cdot\mathbf{x}}}{|\mathbf{k}|^2
      + 2b(T-T_c)} .
  \label{eq:8.14}
\end{equation}
This is just the integral we encountered in our discussion of the Yukawa
potential, Eq.~(4.126).  Evaluating it in the same way, we find
\begin{equation}
  D(\mathbf{x}) = \frac{H_0}{4\pi r}\ e^{-r/\xi},
  \label{eq:8.15}
\end{equation}
where
\begin{equation}
  \xi = \bigl[ 2b(T-T_c) \bigr]^{-1/2}
  \label{eq:8.16}
\end{equation}
is the \emph{correlation length}, the range of correlated spin fluctuations.
Notice that this length diverges as $T\to T_c$.

The main results of this analysis, Eqs.~\eqref{eq:8.5} and \eqref{eq:8.16},
involve unknown constants $b$, $c$ that depend on physics at the atomic scale.
On the other hand, the power-law dependence in these formulae on
$(T-T_c)$ follows simply from the structure of the Landau equations and is
independent of any details of the microscopic physics.  In fact, our
derivation of this dependence did not even use the fact that $G$ describes a
ferromagnet; we assumed only that $G$ can be expanded in powers of an order
parameter and that $G$ respects the reflection symmetry $M\to-M$.  These
assumptions apply equally well to many other types of systems: binary alloys,
superfluids, and even (though the reflection symmetry is less obvious here)
the liquid-gas transition.  Landau theory predicts that the correlation
length behaves as
\begin{equation}
  \xi \sim (T-T_c)^{-\nu},
  \label{eq:8.17}
\end{equation}
with the same exponent $\nu = 1/2$ in all these systems.  This universal
prediction is a simple example of what we expect more generally from the
theory of critical phenomena.

In real physical systems, however, the prediction $\nu = 1/2$ is not quite
correct.  For three-dimensional single-axis magnets, the observed value is
$\nu = 0.630$; for fluids, $\nu = 0.627$.  These values are the same (within
experimental error) for all systems in a given \emph{universality class},
where a universality class is specified by the dimensionality of space, the
number of components of the order parameter, and the behavior of the system
under the reflection symmetry of the order parameter.  For a fluid, the
relevant order parameter is the density, which is a one-component scalar
under reflections, so the fluid belongs to the same universality class as the
single-axis ferromagnet.

Let us try to understand the discrepancy between the Landau prediction and
the observed values of the critical exponents.  Our derivation of
Eq.~\eqref{eq:8.17} neglected the effect of thermal fluctuations.  But exactly
at the critical point, the correlation length diverges, so the fluctuations
are correlated over arbitrarily long distances, and we should not be able to
treat them as a small perturbation.  The proper treatment of these
fluctuations is the subject of the remainder of this book.  Landau theory is
nevertheless a useful starting point, since it correctly predicts the
universality of the critical exponents and the existence of the correlation
length; it only fails to predict the precise numerical values.

To improve on Landau's estimate, we need to take the fluctuations into
account systematically.  The most natural way to do this is to write a
statistical partition function in terms of the continuum field $s(\mathbf{x})$:
\begin{equation}
  Z = \int \mathcal{D}s(\mathbf{x})\ \exp\Bigl[ -\frac{1}{k_BT}\int d^3x\
    \bigl( \tfrac{1}{2}(\nabla s)^2 + b(T-T_c)s^2 + cs^4 - Hs \bigr) \Bigr].
  \label{eq:8.18}
\end{equation}
This is a functional integral over the fluctuating field $s(\mathbf{x})$.  The
appearance of this object signals the deep connection between statistical
mechanics and quantum field theory: The functional integral
\eqref{eq:8.18} is precisely the kind of object that we will use to define
quantum field theories in Chapter~9.  The correspondence is made precise by
identifying the temperature with the inverse imaginary time, as we will
discuss in Section~9.3.

The fluctuations modify the power laws.  For example, Eq.~\eqref{eq:8.16}
is changed, for $T \to T_c$, to
\begin{equation}
  \xi \sim (T-T_c)^{-\nu},
  \label{eq:8.19}
\end{equation}
where the exponent $\nu$ takes a fixed value for all systems in a given
universality class.  For three-dimensional single-axis magnets and for fluids,
$\nu = 0.630$.  The powers in these nontrivial scaling relations are called
\emph{critical exponents}.

The modification from Eq.~\eqref{eq:8.18} to Eq.~\eqref{eq:8.19} does not
imperil the idea that a condensed matter system, in the vicinity of a
second-order phase transition, has a well-defined, cutoff-independent,
continuum behavior.  However, we would like to understand why
Eq.~\eqref{eq:8.17} should be expected as the correct representation.  The
answer to this question will come from a thorough analysis of the ultraviolet
divergences of the corresponding quantum field theory.  In Chapter~12, when we
finally conclude our explication of the ultraviolet divergences, we will find
that we have in hand the tools not only to justify Eq.~\eqref{eq:8.17}, but
also to calculate the values of the critical exponents using Feynman diagrams.
In this way, we will uncover a beautiful application of quantum field theory
to the domain of atomic physics.  The success of this application will guide
us, in Part~III, to even more powerful tools, which we will need in the
relativistic domain of elementary particles.
